% Nouveau EGA V (cote 25): archive cote 25, site page 126 (PDF page 127).
% typescript. Draft compared: none.
% AI-transcribed by Claude Opus 5.5 (claude-opus-5-5, adaptive thinking, effort high), Message Batches API, batch ega25b (msgbatch_01A7n5YuXubUPNK1zsepi8RG), 2026-11-13; not yet checked by eye.
\scanpage{25/126}
\mspage{\hand{29}}

\margin{\unsure{en remarquant, en vérifiant au §22 pour la construction des préliminaires grassmanniennes}}
On peut alors exprimer \over{5}{4}.2. en \unsure{diasant|disant} que l'ensemble des classes
à isomorphisme près de systèmes linéaires de diviseurs sur X associés à
$\underline{L}$ , est en correspondance biunivoque canonique avec l'ensemble \xout{Grassx}
Grass(Q)(S\xout{xxx} et cette correspondance étant compatible avec \unsure{les} changemen
de base quelconques, onvoit que le foncteur
\[ S' \hand{$\rightsquigarrow$} \text{ensemble des classes (mod isomorphisme) de systèmes} \]
\[ \text{linéaires de diviseurs de } X_{S'}/S' \text{ associés à } \underline{L}_{S'} \]
est représentable par le S-préschéma \uline{Grass} (Q).

\margin{\unsure{Sous les conditions de 1.1., on dit que le syst. lin. de diviseurs est complet lorsque $\forall$ $s \in S$, $E(s) \rightarrow H^0(X_s, \underline{L}|X_s)$ est \uline{bijectif}. Ceci ne se conserve pas par changement de base (cf. 4.12) si $g : X \rightarrow S$ est de prés. finie) \& $g$ est propre, alors ceci implique que $g_*(\underline{L})$ est \uline{localement libre de type fini}, et commutation ... $g$ est propre et existence ...}}
\deleted{Le résultat suivant, de nature moins technique que 5.2., devrait
plutôt passer avant ce dernier:

\uline{Corollaire} 5.3. Soit f:X \quad S un morphisme fidèlement plat loc de prés
\inserted[typed]{sur X/S}
finie. Alors pour tout\xout{xxfamille} système linéaire de diviseurs, paramétré
par un fibré projectif P sur S, le morphisme correspondant
\[ \hand{$\check{\ }$}P \qquad \mathrm{Div}_{X/S} \]
est un monomorphisme.

Cela signifie que deux morphismes $S'$ \quad \hand{$\check{\ }$}P ayant même composé avec
le morphisme précédent sont égaux. Faisant le changement de base $S'$ \quad S,
on est ramené \xout{à montrer que deux} au cas de deux sections de \hand{$\check{\ }$}P sur S,
et}

% NOTE: the paragraph from "Il faut expliciter" to "a) Supposons ... algé" is enclosed by hand in a box drawn around it
Il faut expliciter aussi dans \over{5}{4}.1. que $\underline{L}|X-Z$ est canoniquement
isomorphe à $f^{\hand{*}}(\underline{O}_P(1))$ (avec les notations du nº précédent), donc dans
le cas $D \in \emptyset$ envisagé dans \over{5}{4}.3. , $\underline{L}|X-Z$ n'est autre que $\underline{L}_{X-Z}$
de \over{5}{4}.3. b), et par suit e $\underline{L}$ n'est autre que le prolongement canonique et
unique de $f^*(\underline{O}_P(1))$ en un faisceau inversible sur X.

\uline{Proposition} \over{5}{4}.3. Soit D un système linéaire de diviseurs sur X/S
paramétré par \hand{$\check{\ }$}P , où g: X$\longrightarrow$ S est \deleted{\unsure{lisse}} un morphisme plat \inserted{loc.} de présenta-%
tion finie.

a) Supposons \inserted{(g de prés. finie et} que pour tout $s \in S$, désignant par k une clôture algé

\margin{\unsure{lin. complet associé à $\underline{L}$ ssi $g_*(\underline{L})$ est loc. libre, sa formation commutant aux changements} \unsure{16} \unsure{aux fibres, et pour tout $s \in S$, et $\varphi \in H^0(X_s, \underline{L}|X_s)$, $\varphi \neq 0$, $\varphi$ est une section régulière, condi}}
